\chapter{Fazit und Ausblick}

Die vorliegende Arbeit zeigt, dass eine didaktisch wirksame und fachlich belastbare Visualisierung präemptiver CPU-Scheduling-Verfahren als interaktive Webanwendung realisiert werden kann. Die zentrale Leistung besteht nicht nur in der technischen Umsetzung, sondern in der Verbindung von Simulationslogik, erklärender Visualisierung und lernorientierter Bedienung. Dadurch wird die im Betriebssystemkontext zentrale Idee der Prozessverwaltung für Lernende nicht nur theoretisch beschrieben, sondern unmittelbar erfahrbar gemacht.

Die entwickelte Anwendung erfüllt die ursprüngliche Zielsetzung: Sie macht Scheduling-Entscheidungen sichtbar, verknüpft diese mit relevanten Kennzahlen und erlaubt eine nachvollziehbare Gegenüberstellung verschiedener Verfahren auf identischem Szenario. In diesem Sinne ist das Ergebnis ein produktiver Beitrag zur didaktischen Vermittlung von Betriebssystemprinzipien. Die Kombination aus Gantt-Darstellung, Queue-Zustand, Kennzahlen und Vergleichsansicht ist besonders geeignet, um die Dynamik präemptiver Verfahren zu verstehen, ohne die fachliche Genauigkeit zu verlieren.

\section{Ergebnis der Arbeit}

Die Arbeit verbindet theoretische Grundlagen, fachliche Anforderungen, HCI-Überlegungen und technische Umsetzung in einem konsistenten Entwurf. Die Simulation basiert auf einem gemeinsamen, deterministischen Modell, das verschiedene Scheduling-Strategien vergleichbar macht. Gleichzeitig werden die Ergebnisse in einer gestalteten Oberfläche präsentiert, die den Lernprozess durch direkte Visualisierung und nachvollziehbare Zustandswechsel unterstützt.

Das Produkt ist damit nicht bloß eine Demonstrationsumgebung, sondern ein nutzbares Werkzeug für Lehre, Analyse und fachliche Diskussion. Gerade weil die Anwendung unterschiedliche Algorithmen mit derselben Datenbasis vergleicht, wird klar erkennbar, wie sich Entscheidungen in der Auswahlstrategie auf Wartezeit, Durchlaufzeit, Reaktionszeit und Kontextwechsel auswirken.

\section{Ausblick}

Für eine mögliche Weiterentwicklung bietet sich eine Erweiterung in mehreren Richtungen an. Eine erste Erweiterung wäre die Ergänzung weiterer Scheduling-Verfahren oder die Einbindung realistischerer Lastmodelle. Ebenfalls sinnvoll ist die Verfeinerung der Bewertungslogik, etwa durch zusätzliche Gewichtungsoptionen für didaktische oder fachliche Lernziele. Auch eine detailliertere Export- und Analysefunktion könnte die fachliche Nachnutzung der erzeugten Szenarien verbessern.

Ein weiterer sinnvoller Schritt wäre die Überführung der Anwendung in einen stärker didaktisch strukturierten Verlauf, etwa mit erklärenden Lehrpfaden, Aufgabenstellungen oder vorgefertigten Vergleichsszenarien. Dadurch könnte die Anwendung nicht nur als Visualisierungswerkzeug, sondern auch als Lernumgebung mit klarer didaktischer Begleitung fungieren.

Die bestehende Umsetzung bildet daher eine solide Grundlage für weiterführende Arbeiten im Bereich interaktiver Betriebssystemvisualisierung. Sie zeigt, dass technische Korrektheit, Nutzbarkeit und didaktische Klarheit in einem gemeinsamen Produkt zusammengeführt werden können.
